\documentclass[11pt, a4paper]{article}

% --- UNIVERSAL PREAMBLE BLOCK ---
\usepackage[a4paper, top=2.5cm, bottom=2.5cm, left=2cm, right=2cm]{geometry}
\usepackage{fontspec}

\usepackage[english]{babel}

% Set default/Latin font to Sans Serif in the main (rm) slot
\babelfont{rm}{Noto Sans}

% --- Standard Math & Formatting Packages ---
\usepackage{amsmath}
\usepackage{amssymb}
\usepackage{amsthm}
\usepackage{booktabs}
\usepackage{xurl} % Ensures URLs break correctly across lines
\usepackage{graphicx}

% --- Theorem Environments ---
\newtheorem{theorem}{Theorem}[section]
\newtheorem{lemma}[theorem]{Lemma}
\newtheorem{proposition}[theorem]{Proposition}
\newtheorem{corollary}[theorem]{Corollary}
\newtheorem{conjecture}[theorem]{Conjecture}
\theoremstyle{definition}
\newtheorem{definition}[theorem]{Definition}
\newtheorem{remark}[theorem]{Remark}

% Hyperref must be the last package loaded
\usepackage{hyperref}

\title{Computational Bounds and Rational Isomorphisms in the 2-Adic Collatz Flow}
\author{Frank Cortese \\ \small{Adjunct Mathematics Faculty, Shasta Community College}}
\date{\today}

\begin{document}

\maketitle

\begin{abstract}
The Collatz conjecture is typically characterized by non-monotonic, chaotic ``hailstone'' trajectories. We transform this stochastic perspective into a geometric one by introducing a lifted operator $A(n)=3n+2^{v_2(n)}$, which re-indexes the sequence to track the strictly accumulating 2-adic valuation. Expanding the domain to the dyadic ring $\mathbb{Z}[1/2]$ and the rational field $\mathbb{Q}$, we establish a dynamical isomorphism between non-integer rational orbits of the standard map and the integer cycles of generalized $3n+d$ variants. Moving beyond classical Diophantine bounds, we demonstrate that the geometric gap factorizes into an expansion polynomial $P_g$, establishing a strict algebraic isomorphism that maps rational cycles to generalized integer networks. Finally, we introduce the Parity Overhead ($\delta_S$) to quantify the bitwise carry propagation that obstructs the formation of integer cycles, providing a computational metric for the algebraic gaps observed in the integer lattice. In this paper, we present a highly optimized computational architecture for verifying the Collatz Conjecture. Moving beyond standard algebraic bypasses (such as the $4k+3$ heuristic), we introduce an $\mathcal{O}(1)$ Memory Projection Sieve that maps millions of resolved modulo-residue pairs into a contiguous 16 MB array. By strictly sizing this array to fit within modern CPU L3 cache architectures, we achieve a memory-bound look-up latency of approximately one nanosecond. Because this static sieve is pre-computed, it instantly eliminates 98.29\% of the search space before the verification sweep begins. Coupled with Just-In-Time (JIT) compilation to raw machine code, zero-copy shared memory allocation, and dynamic multiprocessing across physical CPU cores, our algorithm avoids standard Inter-Process Communication (IPC) bottlenecks and achieves orders-of-magnitude performance gains over native interpreted loops. We demonstrate the efficacy of this approach by verifying all integers up to $2^{45}$ (approximately 35.1 trillion) on standard consumer hardware, while simultaneously extracting empirical data regarding Maximum Peak Altitudes and Longest Path stopping times.
\end{abstract}

\section{Introduction: Shifting from Probability to Geometry}

The intractability of the Collatz $(3n+1)$ conjecture \cite{Lagarias1985, Lagarias2010} is fundamentally rooted in its location at the intersection of two incompatible arithmetic domains. The iterative process exhibits a dual nature: the multiplicative operation $n \mapsto 3n+1$ represents an Archimedean expansion, driven by the geometric growth of base-3 logic, while the subsequent division by powers of two $n \mapsto n/2^{k}$ constitutes a 2-adic contraction, governed by the valuation metric $|n|_{2}=2^{-v_{2}(n)}$. 

Historically, the interaction between these forces has been modeled as a stochastic process---most notably through the ``hailstone'' metaphor or random walk models. While probabilistic approaches provide an excellent macroscopic description of trajectory decay and the distribution of stopping times, they are inherently limited. By treating the transition between integers as a statistical jump, the rigid algebraic machinery and local topological constraints that dictate the path of a specific orbit are often obscured.

While Bernstein \cite{Bernstein1994} and Lagarias \cite{Lagarias1985} established rigorous $p$-adic and continuous extensions to the Collatz problem, these models often remain separate from the step-by-step integer arithmetic. The introduction of the lifted operator $A(n)=3n+2^{v_{2}(n)}$ bridges this gap by isolating the 2-adic valuation as a strict, additive coordinate. This clarifies exactly what makes our approach structurally novel.

In this paper, we propose a structural shift in perspective. Instead of viewing the trajectory as a series of discrete, non-monotonic jumps in magnitude, we lift the dynamics from the natural numbers $\mathbb{N}$ into the 2-adic integers $\mathbb{Z}_{2}$. This lifting allows us to define a new operator, $A(n) = 3n + 2^{v_2(n)}$. This transformation effectively re-indexes the trajectory to track the strictly accumulating 2-adic valuation.

Our investigation proceeds along two primary avenues:
\begin{enumerate}
    \item \textbf{Topological Extension:} Building upon the $p$-adic extensions explored by Lagarias \cite{Lagarias1985}, Bernstein \cite{Bernstein1994}, and Chamberland \cite{Chamberland1996}, we expand the state space to the Dyadic Ring $\mathbb{Z}[1/2]$ and the rational field $\mathbb{Q}$, revealing that the standard Collatz map naturally extends to the rationals.
    \item \textbf{Dynamical Isomorphisms:} We demonstrate that non-integer rational orbits in the standard map are dynamically isomorphic to the exact integer cycles of generalized $3n+d$ variants. Unlike older Diophantine bounds that merely constrain the existence of cycles, our polynomial expansion mapping proves exact topological equivalencies across distinct arithmetic networks, bridging the discrete and continuous domains.
\end{enumerate}

While large-scale distributed computing projects have extended the brute-force empirical verification of the conjecture to $n = 2^{68}$ via highly optimized bit-shifting runs, such efforts are typically limited to the integer lattice. Our investigation carves a distinct niche in \textbf{High-Precision Rational Manifold Tracking}. Rather than prioritizing raw search-width alone, we focus on the exact preservation of algebraic structures within the dyadic ring $\mathbb{Z}[1/2]$ and the rational field $\mathbb{Q}$, establishing empirical bounds that brute-force records frequently overlook.

In this paper, we present a highly optimized computational architecture for verifying the Collatz Conjecture. Moving beyond standard algebraic bypasses (such as the $4k+3$ heuristic), we introduce an $\mathcal{O}(1)$ Memory Projection Sieve that maps millions of resolved modulo-residue pairs into a contiguous 16 MB array. By strictly sizing this array to fit within modern CPU L3 cache architectures, we achieve a memory-bound look-up latency of approximately one nanosecond. Because this static sieve is pre-computed, it instantly eliminates 98.29\% of the search space before the verification sweep begins. Coupled with Just-In-Time (JIT) compilation to raw machine code, zero-copy shared memory allocation, and dynamic multiprocessing across physical CPU cores, our algorithm avoids standard Inter-Process Communication (IPC) bottlenecks and achieves orders-of-magnitude performance gains over native interpreted loops. We demonstrate the efficacy of this approach by verifying all integers up to $2^{45}$ (approximately 35.1 trillion) on standard consumer hardware, while simultaneously extracting empirical data regarding Maximum Peak Altitudes and Longest Path stopping times.

\section{The Lifted Operator and Strict 2-Adic Monotonicity}

To capture the geometric behavior of these orbits, we must formalize the operator and its relationship to the 2-adic metric. We treat the 2-adic valuation not merely as a division step, but as a fundamental coordinate of the flow.

\begin{definition}[The 2-adic Valuation]
For any positive integer $n\in\mathbb{N}$, let $n=m\cdot2^{k}$ where $m$ is an odd integer and $k\ge0$. The 2-adic valuation is defined as: $v_{2}(n)=k$.
\end{definition}

\begin{definition}[The Lifted Operator on \texorpdfstring{$\mathbb{N}$}{N}]
We define the lifted map $A:\mathbb{N}\rightarrow\mathbb{N}$ as:
\begin{equation}
A(n)=3n+2^{v_{2}(n)}
\end{equation}
\end{definition}

Unlike the standard $3n+1$ step function, $A(n)$ retains the valuation shift as an additive mathematical component. By isolating the valuation shift, we can establish a strict, measurable descent in the 2-adic norm.

\begin{lemma}[Strict 2-adic Monotonicity]
\label{lem:monotone}
The 2-adic norm $|n|_{2}=2^{-v_{2}(n)}$ is a strict Lyapunov function for $A(n)$ on the domain of positive integers excluding the invariant set $T = \{2^k \mid k \in \mathbb{Z}_{\ge 0}\}$.
\end{lemma}

\begin{proof}
Let $n \in \mathbb{N} \setminus T$. Then $n = m \cdot 2^{k}$ for some odd $m > 1$. Applying the operator yields $A(n)=2^{k}(3m+1)$. Since $m>1$ is odd, $3m+1$ is an even integer greater than 4, thus $v_{2}(3m+1)\ge1$. The new altitude is $v_{2}(A(n))=k+v_{2}(3m+1)\ge k+1$. In terms of the 2-adic metric, $|A(n)|_{2} < |n|_{2}$.
\end{proof}

\begin{definition}[The Odd Kernel]
For any $n\in\mathbb{N}$, the odd kernel $\phi(n)$ is obtained by stripping away the 2-adic valuation:
$\phi(n)=n/2^{v_{2}(n)}$.
\end{definition}

\begin{definition}[The Syracuse Map]
The Syracuse map $T_{\text{Syr}} : \mathbb{O}^{+} \rightarrow \mathbb{O}^{+}$ is defined as:
$T_{\text{Syr}}(m) = \phi(3m+1) = (3m+1)/2^{v_2(3m+1)}$.
This map represents the standard Collatz dynamics restricted to the odd lattice.
\end{definition}

\begin{theorem}[Dynamical Semi-Conjugacy]
\label{thm:semi_conjugacy}
The operator $A(n)$ is semi-conjugate to the Syracuse map $T_{\text{Syr}}$ via the projection $\phi$.
\end{theorem}

The semi-conjugacy established in Theorem \ref{thm:semi_conjugacy} implies that the lifted operator preserves the underlying Syracuse dynamics while re-indexing the flow to track the accumulating valuation. This relationship is summarized by the commutative diagram in Figure \ref{fig:semi_conjugacy}, where the path $T_{\text{Syr}} \circ \phi$ is equivalent to $\phi \circ A$, effectively projecting the strictly monotonic geometric flow back onto the classical odd lattice.

\begin{figure}[ht]
    \centering
    \[
    \begin{array}{ccc}
    \mathbb{N} & \xrightarrow{A} & \mathbb{N} \\
    \downarrow{\phi} & & \downarrow{\phi} \\
    \mathbb{O}^{+} & \xrightarrow{T_{\text{Syr}}} & \mathbb{O}^{+}
    \end{array}
    \]
    \caption{Commutative diagram illustrating the semi-conjugacy between $A$ and $T_{\text{Syr}}$ via the projection \(\phi\).}
    \label{fig:semi_conjugacy}
\end{figure}

\section{Expanding the Domain: The Dyadic Ring and Primitive Sets}

We expand the state space to the Dyadic Ring, $\mathbb{Z}[1/2] = \{a/2^k \mid a \in \mathbb{Z}, k \in \mathbb{Z}_{\ge 0}\}$.

\begin{definition}[Primitive Dyadic Sets]
We define the set of Primitive Dyadic Rationals at precision $S$ as:
$\mathcal{P}_{S}=\{a/2^{S} \mid a \in \mathbb{O}^{+}, a < 2^{S}\}$.
\end{definition}

\begin{proposition}[Dynamical Isomorphism]
\label{thm:isomorphism}
The union $\mathcal{U} = \bigcup \mathcal{P}_S$ is dynamically isomorphic to $\mathbb{O}^{+}$. For any $n\in\mathbb{O}^{+}$ the trajectory of $n$ under $A(n)$ is identical to that of its representative $n/2^{S} \in \mathcal{P}_{S}$ up to a scaling factor $2^{-S}$.
\end{proposition}

\begin{proof}
Let $n \in \mathbb{O}^+$ be an arbitrary odd integer, and let $q = n/2^S \in \mathcal{P}_S$ be its representation in the Primitive Dyadic set for some precision $S \in \mathbb{Z}_{\ge 0}$. We examine the application of the lifted operator $A$ to the dyadic rational $q$. By the properties of the 2-adic valuation, $v_2(q) = v_2(n/2^S) = v_2(n) - S$. Applying the operator yields:
\begin{equation*}
A(q) = 3\left(\frac{n}{2^S}\right) + 2^{v_2(n/2^S)} = \frac{3n}{2^S} + 2^{v_2(n)-S} = \frac{3n + 2^{v_2(n)}}{2^S} = \frac{A(n)}{2^S}
\end{equation*}
Algebraically, the 2-adic valuation operator $v_2$ operates identically on the numerator $n$, meaning the sequence of operations perfectly factors out the $2^{-S}$ term at every iterative step. By induction, this scalar relationship holds for all subsequent iterations. Thus, the trajectory of $q$ maintains exact dynamic parity with the integer trajectory of $n$, establishing the isomorphism.
\end{proof}

\section{The Rational Extension and the \texorpdfstring{$3n+d$}{3n+d} Isomorphism}

\begin{definition}[The Generalized $3n+d$ Operator]
For any odd integer $d \in \mathbb{Z}^{+}$ the generalized Collatz operator $T_d : \mathbb{Z} \rightarrow \mathbb{Z}$ is defined as:
$T_d(n) = 3n + d \cdot 2^{v_2(n)}$.
\end{definition}

\begin{lemma}[Universal Invariance of the Denominator]
\label{lem:denominator}
Any positive rational $q \in \mathbb{Q}^{+}$ can be represented as $q = a/d$. Let $d = d_{\text{odd}} \cdot 2^{k}$. The iteration of $A(q)$ preserves the denominator $d$, while the numerator evolve according to the operator $T_{d_{\text{odd}}}$.
\end{lemma}

\begin{table}[htbp]
\centering
\caption{Full Trace of Rational Orbits under $A(q)$ Mirroring Integer $T_d$ Dynamics}
\label{tab:isomorphism_trace}
\begin{tabular}{l l l l l}
\toprule
Step ($i$) & Rational State $q_i$ ($3n+1$) & Integer Num. $a_i$ (Target $T_d$) & Odd Kernel $\phi(q_i)$ & $d_{\text{odd}}$ \\ \midrule
\multicolumn{5}{c}{\textit{Family: $d_{\text{odd}} = 7$}} \\
0 & $5/7$ & 5 & $5/7$ & 7 \\
1 & $22/7$ & 22 & $11/7$ & 7 \\
2 & $80/7$ & 80 & $\mathbf{5/7}$ (Cycle) & 7 \\ \midrule
\multicolumn{5}{c}{\textit{Family: $d_{\text{odd}} = 65$}} \\
0 & $8/65$ & 8 & $1/65$ & 65 \\
1 & $544/65$ & 544 & $17/65$ & 65 \\
2 & $3712/65$ & 3712 & $29/65$ & 65 \\
3 & $19456/65$ & 19456 & $19/65$ & 65 \\
4 & $124928/65$ & 124928 & $61/65$ & 65 \\
5 & $507904/65$ & 507904 & $31/65$ & 65 \\
6 & $2588672/65$ & 2588672 & $79/65$ & 65 \\
7 & $9895936/65$ & 9895936 & $151/65$ & 65 \\
8 & $33947648/65$ & 33947648 & $259/65$ & 65 \\
9 & $110362624/65$ & 110362624 & $421/65$ & 65 \\
10 & $348127232/65$ & 348127232 & $83/65$ & 65 \\
11 & $1317011456/65$ & 1317011456 & $157/65$ & 65 \\
12 & $4496293888/65$ & 4496293888 & $67/65$ & 65 \\
13 & $17850957824/65$ & 17850957824 & $133/65$ & 65 \\
14 & $62277025792/65$ & 62277025792 & $\mathbf{29/65}$ (Cycle) & 65 \\ \midrule
\multicolumn{5}{c}{\textit{Family: $d_{\text{odd}} = 1$}} \\
0 & $7/8$ & 7 & $7$ & 1 \\
1 & $11/4$ & 11 & $11$ & 1 \\
2 & $17/2$ & 17 & $17$ & 1 \\
3 & $26$ & 26 & $13$ & 1 \\
4 & $80$ & 80 & $5$ & 1 \\
5 & $256$ & 256 & $\mathbf{1}$ (Collapsed) & 1 \\
\bottomrule
\end{tabular}
\end{table}

\begin{proposition}[The \texorpdfstring{$3n+d$}{3n+d} Dynamical Isomorphism]
\label{thm:general_iso}
The trajectory of any $q = a/d \in \mathbb{Q}$ under $A$ is identical to the scaled trajectory of the integer sequence generated by the $T_{d_{\text{odd}}}$ map starting from the integer seed $a$.
\end{proposition}

\begin{proof}
Let $q_i = a_i/d$ be an arbitrary rational state within a trajectory, where $d = d_{\text{odd}} \cdot 2^k$. As established in Lemma \ref{lem:denominator}, the 2-adic valuation is $v_2(q_i) = v_2(a_i) - v_2(d) = v_2(a_i) - k$. Applying the lifted operator $A$ to $q_i$ gives:
\begin{equation*}
A(q_i) = 3\left(\frac{a_i}{d}\right) + 2^{v_2(a_i)-k} = \frac{3a_i + d \cdot 2^{v_2(a_i)-k}}{d}
\end{equation*}
Substituting $d = d_{\text{odd}} \cdot 2^k$ into the additive term in the numerator, we obtain:
\begin{equation*}
A(q_i) = \frac{3a_i + (d_{\text{odd}} \cdot 2^k) \cdot 2^{v_2(a_i)-k}}{d} = \frac{3a_i + d_{\text{odd}} \cdot 2^{v_2(a_i)}}{d}
\end{equation*}
We observe that multiplying the resulting rational orbit state by the denominator yields the exact formulation of the generalized operator: $d \cdot A(q_i) = 3a_i + d_{\text{odd}} \cdot 2^{v_2(a_i)} = T_{d_{\text{odd}}}(a_i)$. Because the operator commutes with this scalar multiplication, the rational trajectory $\{q_i\}$ and the integer sequence $\{a_i\}$ generated by $T_{d_{\text{odd}}}$ are strictly isomorphic. Therefore, a cycle of length $L$ exists in $\mathbb{Q}$ if and only if a corresponding integer cycle of length $L$ exists for $T_{d_{\text{odd}}}$.
\end{proof}

\section{Quantifying the Obstruction: Parity Overhead}

To exist, an integer cycle of length $L$ and total valuation shift $S$ must satisfy the global cycle equation \cite{Steiner1977}:
\begin{equation}
\label{eq:cycle_raw}
n_0 = \frac{\sum_{j=0}^{L-1} 3^{L-1-j} 2^{k_j}}{2^S - 3^L} = \frac{\mathcal{N}}{\mathcal{G}}
\end{equation}

\subsection{Structural Source: Bitwise Carry Propagation}
In binary arithmetic, the addition operation in the $3n+1$ step frequently generates bitwise carries. These carries propagate through the binary representation, unexpectedly inflating the 2-adic valuation shift required to return to an odd integer. We define the \textbf{Diophantine Minimum} as the smallest integer $S$ that satisfies $2^S > 3^L$: $S_{min} = \lceil L \log_2 3 \rceil$. Any observed trajectory resulting in a shift $S_{obs}$ exceeding this is said to have a \textbf{Parity Overhead} ($\delta_S = S_{obs} - S_{min}$). In this regime, the system is governed by bounds on linear forms in logarithms \cite{Matveev2000}.

According to Baker's theory of linear forms in logarithms (specifically Matveev's theorem), the logarithmic distance between powers of 2 and 3 is strictly bounded from below by an exponential function of the cycle length $L$:
\begin{equation}
|S_{obs} \ln 2 - L \ln 3| > \exp(-C \ln L)
\end{equation}
for some computable absolute constant $C > 0$. By algebraically substituting our definition of Parity Overhead ($S_{obs} = S_{min} + \delta_S$) into this bound, we explicitly isolate the algebraic obstruction:
\begin{equation}
|(\delta_S + S_{min}) \ln 2 - L \ln 3| = |\delta_S \ln 2 + (S_{min} \ln 2 - L \ln 3)| > \exp(-C \ln L)
\end{equation}
For an integer cycle to exist, the geometric gap $\mathcal{G} = 2^S - 3^L$ must remain small enough to divide the bounded algebraic noise $\mathcal{N}$, necessitating a logarithmic distance near zero. The component $(S_{min} \ln 2 - L \ln 3)$ represents the baseline irrationality measure between the base-2 and base-3 expansions. Because the theoretical lower bound restricts how tightly this baseline can approximate zero for large values of $L$, introducing any positive Parity Overhead ($\delta_S \ge 1$) drastically inflates the overall magnitude by discrete leaps of $\ln 2 \approx 0.693$. This demonstrates precisely how the theoretical lower bound severely restricts the existence of $\delta_{S}$ for large values of $L$: a positive overhead creates a severe algebraic obstruction, forcing the geometric gap to exponentially outgrow the expected algebraic noise. This severely restricts the solution interval, acting as a primary deterministic bound on the formation of integer cycles for that specific seed.

\subsection{The Interval Collapse Mechanism}
The Parity Overhead serves as a \textbf{Path-Dependent Algebraic Obstruction}, quantifying the exclusion of specific integer paths from cyclic orbits. We define the \textbf{Solution Interval} ($\Delta n$) as a measure of \textbf{Solution Density}, representing the normalized bounds window available for a cycle of given length $L$ and shift $S$:
\begin{equation}
\Delta n = \frac{3^L}{2^S - 3^L}
\end{equation}
As $\delta_S$ increases, the denominator $2^S - 3^L$ grows exponentially. For non-cyclic integer seeds (such as $n=27$ or $31$), we evaluate this hypothetical window using their observed trajectory lengths and shifts. Because the overhead inflates the geometric gap, the density "collapses" to a width significantly less than unity ($\Delta n \ll 1$). Since the path-dependent noise $\mathcal{N}$ in the cycle equation is fundamentally bounded by the sequence's geometric growth, this structural collapse acts as a profound deterministic bound on the formation of a closed integer loop for any sequence where $\delta_S \ge 1$, providing a strict structural limit that precludes periodicity and drives the empirically observed trajectory decay along these specific paths.

\begin{table}[htbp]
\centering
\caption{Parity Overhead and Solution Interval Density for Selected Trajectories}
\label{tab:interval_collapse}
\begin{tabular}{l c c c c c l}
\toprule
Seed ($n$) & $L$ & $S_{min}$ & $S_{obs}$ & Overhead ($\delta_S$) & Density ($\Delta n$) & Classification \\ \midrule
\multicolumn{7}{c}{\textit{Rational Control Group}} \\
$1/3$ & 1 & 2 & 2 & 0 & $3.000 \times 10^{0}$ & Rational Cycle \\
$13/11$ & 1 & 2 & 2 & 0 & $3.000 \times 10^{0}$ & Rational Cycle \\
$19/5$ & 3 & 5 & 5 & 0 & $5.400 \times 10^{0}$ & Rational Cycle \\ 
$31/32$ & 39 & 62 & 67 & 5 & $2.824 \times 10^{-2}$ & Observed Decay ($\Delta n \ll 1$) \\ \midrule
\multicolumn{7}{c}{\textit{Integer Lattice Targets}} \\
3 & 5 & 8 & 8 & 0 & $1.869 \times 10^{1}$ & Potential Loop \\
7 & 11 & 18 & 20 & 2 & $2.033 \times 10^{-1}$ & Observed Decay ($\Delta n \ll 1$) \\
27 & 41 & 65 & 70 & 5 & $3.188 \times 10^{-2}$ & Observed Decay ($\Delta n \ll 1$) \\
31 & 39 & 62 & 67 & 5 & $2.824 \times 10^{-2}$ & Observed Decay ($\Delta n \ll 1$) \\
31948311 & 111 & 176 & 189 & 13 & $1.164 \times 10^{-4}$ & Observed Decay ($\Delta n \ll 1$) \\
$2^{45} - 1$ & 386 & 612 & 615 & 3 & $1.217 \times 10^{-1}$ & Observed Decay ($\Delta n \ll 1$) \\
\bottomrule
\end{tabular}
\end{table}

\section{Algebraic Factorization of the Geometric Gap}

Let $g = \gcd(L, S) \ge 1$. The gap factorizes into polynomial components:
\begin{equation}
\mathcal{G} = 2^{gs} - 3^{g\ell} = (2^s - 3^\ell) \sum_{i=0}^{g-1} (2^s)^{g-1-i} (3^\ell)^i = \mathcal{G}_{prim} \cdot P_g
\end{equation}

The polynomial expansion $P_g$ derived from the gap $\mathcal{G}$ is not just an algebraic artifact, but acts as the precise scaling factor. This factor maps the rational cyclic domain directly into the generalized integer operator $T_{P_g}$. Table \ref{tab:target_mapping} explicitly maps these resonant lengths to their corresponding polynomial expansions to illustrate this relationship.

\begin{table}[htbp]
\centering
\caption{Target Mapping: Resonant Lengths and their Expansion Polynomials ($P_g$)}
\label{tab:target_mapping}
\begin{tabular}{cccllr}
\toprule
$L$ & $S$ & $g$ & Gap \(\mathcal{G} = 2^S - 3^L\) & Factorization (\(\mathcal{G}_{prim} \times P_g\)) & Target Operator \\ \midrule
1   & 3   & 1   & 5                               & $5 \times 1$                                      & $T_{5}$         \\
2   & 4   & 2   & 7                               & $1 \times 7$                                      & $T_{7}$         \\
3   & 5   & 1   & 5                               & $5 \times 1$                                      & $T_{5}$         \\
6   & 10  & 2   & 295                             & $5 \times 59$                                     & $T_{59}$        \\
9   & 15  & 3   & 13,085                          & $5 \times 2,617$                                  & $T_{2,617}$     \\
10  & 16  & 2   & 6,487                           & $13 \times 499$                                   & $T_{499}$       \\
15  & 25  & 5   & 19,205,525                      & $5 \times 3,841,105$                              & $T_{3,841,105}$ \\
\bottomrule
\end{tabular}
\end{table}

\begin{theorem}[Algebraic Isomorphism of Rational Orbits]
\label{thm:pg_iso}
Let $q = a/P_g$ be a rational cycle of the standard $3n+1$ operator $A$. The mapping \(\psi(q) = P_g \cdot q\) is a dynamical isomorphism onto an integer cycle of $T_{P_g}$.
\end{theorem}

\begin{proof}
Recall the factorization of the geometric gap $\mathcal{G} = \mathcal{G}_{prim} \cdot P_g$. Let $q$ be a rational cycle of the standard $3n+1$ operator $A$ possessing the Expansion Polynomial $P_g$ as its denominator, such that $q_i = a_i / P_g$ for $a_i \in \mathbb{Z}$. Applying the mapping \(\psi(q_i) = P_g \cdot q_i = a_i\) completely clears the denominator. Because $P_g$ is derived as an odd integer (consisting of sums of powers of 3 and 2), its 2-adic valuation is zero ($v_2(P_g) = 0$).

Using the logic established in the proof for Proposition \ref{thm:general_iso}, the valuation of the rational state is exactly $v_2(q_i) = v_2(a_i) - v_2(P_g) = v_2(a_i)$. Applying the standard map $A$:
\begin{equation*}
A(q_i) = 3\left(\frac{a_i}{P_g}\right) + 2^{v_2(a_i)} = \frac{3a_i + P_g \cdot 2^{v_2(a_i)}}{P_g}
\end{equation*}
The numerator is precisely the definition of the generalized operator $T_{P_g}(a_i)$. Consequently, $A(q_i) = T_{P_g}(a_i) / P_g$. This demonstrates that the scalar multiplication by $P_g$ mathematically forces the numerator into an exact integer cycle within the $T_{P_g}$ domain, cementing the algebraic isomorphism. \qed
\end{proof}

\section{Computational Methodology and Verification}

To bridge the gap between theoretical isomorphisms and empirical observation, we developed a custom, fully parallelized Python-based computational suite. While world-record distributed searches reaching $n=2^{68}$ utilize low-level bit-shifting to maximize throughput, our architecture is optimized for absolute numerical fidelity across the rational extension. We achieve continuous verification up to $n=2^{45}$ while maintaining $S=10,000$ precision for the mapping of rational manifolds, preserving the denominator $d_{\text{odd}}$ as an exact mathematical entity throughout the flow---a requirement absent in standard brute-force records. Unlike integer-only distributed sieves, our methodology necessitates the exact preservation of the denominator $d_{odd}$ across $S=10,000$ iterations to map the local topological curvature of the rational manifold.

\subsection{Data Structures and Exact 2-Adic Precision}
The standard floating-point architecture (IEEE 754) is fundamentally incompatible with the precision required to track 2-adic flows across thousands of iterations, as intermediate values quickly exceed the 53-bit mantissa limit.

To prevent catastrophic precision loss, the state space of the dyadic ring $\mathbb{Z}[1/2]$ and the rational field $\mathbb{Q}$ is represented using Python's native arbitrary-precision integers and the \texttt{fractions.Fraction} module\footnote{Python Software Foundation, \texttt{fractions} --- Rational numbers: \url{https://docs.python.org/3/library/fractions.html}}. This guarantees that the denominator $d_{\text{odd}}$ is preserved as an exact mathematical entity throughout the iteration.

The core computational bottleneck---calculating the 2-adic valuation $v_{2}(n)$---is optimized by bypassing string conversions or modulo arithmetic. Instead, the valuation is extracted in $\mathcal{O}(1)$ time using native two's complement bitwise logic to isolate the least significant set bit: \texttt{(n \& -n).bit\_length() - 1}\footnote{Python Software Foundation, Built-in Types (\texttt{int.bit\_length}): \texttt{https://docs.python.org/3/library/stdtypes.html\#int.bit\_length}}. For dyadic rationals, this is extended algebraically as $v_2(\text{numerator}) - v_2(\text{denominator})$.

\subsection{Algorithmic Complexity and Trajectory Tracking}
Applying the lifted operator $A(q) = 3q + 2^{v_{2}(q)}$ requires tracking the ``Odd Kernel'' $\phi(q_i)$ at each step to detect cyclic behavior or interval collapse.

\textbf{Cycle Detection:} For rational orbits (as demonstrated in Table \ref{tab:isomorphism_trace}), the algorithm stores previous odd kernels in an array. A state is evaluated as strictly cyclic if a newly generated \(\phi(q_i)\) matches a historically visited kernel, or as ``collapsed'' if \(\phi(q_i) = 1\).

\textbf{Bounded Parity Overhead:} To evaluate the Solution Interval $(\Delta n)$ for large integer seeds, the algorithm isolates the Diophantine Minimum $S_{min} = \lceil L \log_2 3 \rceil$ using the \texttt{math} library\footnote{Python Software Foundation, \texttt{math} --- Mathematical functions: \url{https://docs.python.org/3/library/math.html}}. The Parity Overhead $(\delta_S)$ and the geometric gap $\mathcal{G} = |2^{S_{obs}} - 3^L|$ are derived explicitly to definitively classify trajectory decay (where $\Delta n \ll 1$) versus cyclic potential.

\textbf{Large-Scale Integer Sieving:} Validating the integer lattice up to $n=2^{45}$ via brute-force linear progression is computationally unfeasible. To optimize this, the algorithm utilizes a mathematically constrained generator. Because any odd seed of the form $n \equiv 1 \pmod 4$ strictly drops below its initial value after two 2-adic descent steps ($3(4k+1)+1 = 12k+4 \rightarrow 3k+1 < 4k+1$), these integers are structurally guaranteed to decay. The algorithm's search space is therefore initialized at $n \equiv 3 \pmod 4$ and configured to step by 4, completely bypassing the generation and evaluation of $4k+1$ seeds. For the remaining highly volatile $4k+3$ seeds, the algorithm relies on the ``drop-below'' heuristic. The evaluation of a massive seed terminates the exact moment its trajectory drops below its initial starting value, confirming intersection with a previously verified bounded path.

\subsection{The \texorpdfstring{$\mathcal{O}(1)$}{O(1)} Memory Projection Sieve and Cache Optimization}

Traditional sieving methods for the Collatz Conjecture rely on iterative mathematical checks that incur significant CPU overhead. To maximize throughput, we shift this computational burden to constant-time hardware-level memory retrieval. We pre-computed a breadth-first expansion of the Collatz tree to identify ``resolved'' congruence classes—integer branches mathematically guaranteed to drop below their initial seed within a finite number of steps. While $n \equiv 1 \pmod{4}$ is a trivial case, we extended this analysis to higher-order residues such as $16k+3$ and $16k+11$, up to a modulo depth of $2^{24}$.

These millions of resolved residue pairs are projected into a single, flat, one-dimensional boolean array of size $2^{24}$ (16 MB), allocated directly at the operating system level using zero-copy shared memory. This array size was chosen specifically to align with the L3 cache specifications of modern multi-core processors. By ensuring the entire lookup table resides within the processor's high-speed SRAM, the algorithm avoids the severe ``latency penalty'' (typically $>50$ ns) associated with fetching data from main system RAM (DRAM). When testing an integer $n$, the algorithm performs a single hardware-bound $\mathcal{O}(1)$ retrieval: 
\begin{equation}
\text{is\_resolved} = \text{array}[n \pmod{2^{24}}]
\end{equation}
Because this array is backed by OS-level shared RAM, an arbitrary number of parallel worker processes can read the exact same physical memory block simultaneously without incurring serialization overhead or duplicating the data across cores. This L3-bound projection allows for a lookup latency of approximately 1 ns. Because the sieve is a static, pre-computed structure, it effectively bypasses the von Neumann bottleneck and instantly eliminates 98.29\% of the $2^{45}$ search space before the sweep begins, without requiring a single arithmetic simulation for the bypassed seeds.

\subsection{Just-In-Time (JIT) Compilation and 2-Adic Extraction}

For the remaining 1.71\% of the search space---classified as ``Indeterminate'' seeds by the L3 Cache Sieve---the integers were subjected to explicit Collatz simulation. To bypass the performance bottleneck inherent to Python's interpreted execution, the core verification loop was wrapped in a Just-In-Time (JIT) compiler via the LLVM-based \texttt{numba} library \cite{Lam2015}.

By applying the \texttt{@njit} decorator, the Python syntax was translated directly into highly optimized native machine code prior to execution. To maximize throughput, we replaced high-level object methods (such as \texttt{.bit\_length()}) with continuous bitwise right-shifts (\texttt{while (current \& 1) == 0: current >>= 1}). During compilation, the LLVM backend aggressively optimizes this bitwise loop into single hardware-level machine instructions---specifically \texttt{TZCNT} (Trailing Zero Count) or \texttt{BSF} (Bit Scan Forward) on x86 architectures. This exact hardware-level translation ensures highly efficient, constant-time 2-adic extraction. Further algorithmic efficiency was achieved by implementing a mathematical fused step within the inner loop. Because any odd integer $n$ guarantees that $3n+1$ is strictly even, the subsequent 2-adic division step is deterministic. By fusing the multiplication and the resulting bitwise right-shift into a single continuous operation, we eliminate a redundant branch evaluation on every odd step. Finally, because our modulo-indexing is mathematically constrained, we passed the \texttt{boundscheck=False} directive to the JIT compiler\footnote{Numba Documentation, Performance Tips --- No bounds checking: \texttt{https://numba.pydata.org/numba-doc/latest/user/performance-tips.html\#no-bounds-checking}}, allowing the LLVM backend to safely strip away redundant array-boundary safety checks during the high-speed L3 cache lookups.

Furthermore, by passing the \texttt{nogil=True} parameter to the JIT compiler, we explicitly released Python's Global Interpreter Lock (GIL). Bypassing the GIL is critical; when combined with Python's \texttt{multiprocessing.Pool} architecture\footnote{Python Software Foundation, \texttt{multiprocessing} --- Process-based parallelism: \url{https://docs.python.org/3/library/multiprocessing.html}}, it enabled an embarrassingly parallel workload distribution. This dynamic chunking system automatically scales to the host hardware, allowing the verification suite to achieve a continuous 100\% CPU utilization across all available physical cores when evaluating the indeterminate seeds.

\subsection{Execution Environment}
The entire computational verification pipeline was executed within a Google Colab Pro environment. This infrastructure provided the extended RAM necessary to handle the memory-intensive caching of massive trajectory histories and allowed for accelerated runtime execution without local hardware limitations. Because the \texttt{multiprocessing.Pool} implementation dynamically polls the host hardware, the architecture automatically scaled to utilize the 8 to 12 vCPUs provisioned by the Colab Pro High-RAM instances, seamlessly distributing the workload without requiring hardcoded thread limits.

\subsection{Open Source and Institutional Scalability}
The continuous verification of the integer lattice up to $2^{45}$ on consumer hardware serves primarily as a baseline proof-of-concept for our computational architecture. By making the complete verification suite and its underlying methodology openly available via a public GitHub repository, we aim to provide the mathematical community with a transparent, fully reproducible toolset for high-precision rational manifold tracking.

Because the $\mathcal{O}(1)$ Memory Projection Sieve is inherently architecture-agnostic and specifically designed around universal L3 cache mapping principles, the software is perfectly primed for deployment across high-performance university computing clusters. Unburdened by the von Neumann bottleneck and coupled with our JIT-compiled extraction loops, scaling this parallelized architecture across institutional supercomputers will allow researchers to effortlessly push empirical verification well past current computational bounds.

\section{Results and Empirical Discoveries}

The implementation of the L3-bound $\mathcal{O}(1)$ Sieve and JIT compilation yielded transformative performance improvements. A baseline verification of $2^{35}$ natively required extensive compute time; our optimized architecture verified the same range in 1.43 seconds---an approximate 600x speedup.

Pushing the architecture to its limits, we conducted a continuous verification of the search space up to $2^{45}$ (35,184,372,088,832). By distributing the 35.1 trillion integers into manageable, asynchronous 1-billion integer chunks across the CPU core pool, the parallelized sweep completed the verification in 6 hours, 16 minutes, and 12.09 seconds.

During this sweep, we successfully extracted empirical bounding data for the entire domain. The maximum peak altitude encountered in the $2^{45}$ range was 9,223,371,995,968,455,082, which originated from the initial seed 11,390,726,016,927. Furthermore, the longest sequence required to drop below the initial starting value consisted of 949 steps, originating from the seed 14,616,588,676,251. These empirical bounds reinforce the behavioral consistency of the Collatz trajectories at the trillion-scale level and provide concrete targets for future heuristic analysis.

\section{Conclusion}

The transition from a stochastic to a geometric perspective of the Collatz conjecture reveals that the traditionally chaotic ``hailstone'' trajectories are shadows of a higher-dimensional monotonic flow. By lifting the map into the 2-adic integers via the operator $A(n) = 3n + 2^{v_2(n)}$, we have demonstrated that the underlying arithmetic is governed by a strict, measurable descent in the 2-adic norm. 

As visually demonstrated in Figure \ref{fig:hailstone_vs_monotonic}, this re-indexing strips away the Archimedean noise of the standard hailstone jumps. The volatile fluctuations of seeds like $n=27$ are transformed into a strictly increasing trajectory, directly visualizing the continuous accumulation of the 2-adic valuation. This ``Monotonic Increasing'' behavior is not merely a formal curiosity but a structural bridge; we have established through the $3n+d$ Dynamical Isomorphism that non-integer rational orbits in the standard map are the exact structural equivalents of the integer cycles found in generalized $T_{P_g}$ variants. 

Ultimately, the existence---or absence---of integer cycles is determined by the Parity Overhead ($\delta_S$). This metric quantifies the bitwise carry propagation that separates the integer lattice from the periodic orbits found in the rational extension. As evidenced by the Solution Interval ($\Delta n$) collapse in Table \ref{tab:interval_collapse}, a positive overhead forces a geometric gap that exponentially outgrows the expected algebraic noise. This creates a severe algebraic obstruction, acting as a primary topological limit against the formation of closed integer loops for any trajectory exhibiting this positive overhead. 

Future experimental work should focus on the distribution of $\delta_S$ across larger computational manifolds to further refine the deterministic bounds on the linear forms in logarithms that isolate the integer lattice.

\begin{figure}[htbp]
    \centering
    \IfFileExists{Figure_2_Publication_Ready.pdf}{
        \includegraphics[width=0.9\textwidth]{Figure_2_Publication_Ready.pdf}
    }{
        \framebox{\parbox{0.9\textwidth}{\centering
            \vspace{3cm}
            \textbf{MISSING FIGURE: Figure\_2\_Publication\_Ready.pdf} \\
            \small\textit{Side-by-side Analysis: Hailstone vs. Monotonic Increasing\\(Generated via Python Suite for Seed 27)}
            \vspace{3cm}
        }}
    }
    \caption{Side-by-side trajectory analysis of seed $n=27$. (Left) The State Hailstone exhibiting non-monotonic fluctuations in the standard Syracuse map. (Right) The Monotonic Increasing trajectory of the Uplift Operator $A(n) = 3n + 2^{v_2(n)}$ demonstrating strict 2-adic monotonicity.}
    \label{fig:hailstone_vs_monotonic}
\end{figure}

\begin{thebibliography}{9}
\bibitem{Bernstein1994} D. J. Bernstein, \textit{A non-iterative 2-adic statement of the 3N+1 conjecture}, Proceedings of the American Mathematical Society, Vol. 121, No. 2, pp. 405--408, 1994.
\bibitem{Chamberland1996} M. Chamberland, \textit{A continuous extension of the 3x+1 problem to the real line}, Dynamics of Continuous, Discrete and Impulsive Systems, Vol. 2, No. 4, pp. 495--509, 1996.
\bibitem{Lagarias1985} J. C. Lagarias, \textit{The 3x + 1 problem and its generalizations}, The American Mathematical Monthly, Vol. 92, No. 1, pp. 3--23, 1985.
\bibitem{Lagarias2010} J. C. Lagarias (Ed.), \textit{The Ultimate Challenge: The 3x+1 Problem}, 2010.
\bibitem{Lam2015} S. K. Lam, A. Pitrou, and S. Seibert, \textit{Numba: A LLVM-based Python JIT compiler}, Proceedings of the Second Workshop on the LLVM Compiler Infrastructure in HPC, pp. 1--6, 2015.
\bibitem{Matveev2000} E. M. Matveev, \textit{Explicit lower bounds for linear forms in logarithms}, Izvestiya: Mathematics, Vol. 64, No. 6, pp. 1217--1269, 2000.
\bibitem{Steiner1977} R. P. Steiner, \textit{A theorem on the syracuse problem}, Proceedings of the 8th Manitoba Conference on Numerical Mathematics, pp. 553--559, 1977.
\end{thebibliography}

\end{document}
